\section{A finite colored-graph characterization}\label{sec:graph}

We first give a criterion whose individual witnesses can be checked by finite
combinatorial calculations. Throughout this section, $L$ is a finite nonempty
lattice, with least element $0$.

The graph language follows the valuation and path approach of
Pudl\'ak \citep[Definitions~1--3 and pp.~271--272]{Pudlak1976}; the
unary-closure criterion is recorded in
\citet[Theorem~3.2.1]{DeMeo2012}.  We give the full finite argument,
including its witness-size consequence.  Kenney's graphical-algebra
formulation contains the label-decreasing-map and edge-image-path
reconstruction used here
\citep{Kenney2010}\citep[pp.~51--58]{Kenney2009Slides}.
Our statement isolates a finite test without an auxiliary graphical
composition operation.

\begin{definition}[Admissible placements]\label{f:def-placement}
An $L$-colored complete graph consists of a finite nonempty vertex set $A$ and a
symmetric function $d:A\times A\to L$, with $d(p,p)=0$. Loops are included among
its edges. A placement into another such graph $(A',d')$ is a map $f:A\to A'$;
it is \emph{admissible} if
\[
 d'(f(p),f(q))\le d(p,q)\qquad(p,q\in A).
\]
For a nonempty set $S$ of edges of $(A,d)$, let $\Gamma_S$ be the undirected
graph on $A$ whose marked adjacencies are the pairs
$\{f(p),f(q)\}$, where $\{p,q\}\in S$ and $f:A\to A$ is admissible.
Collapsed pairs are allowed and do not affect connectivity. Write
$a_S=\bigvee_{\{p,q\}\in S}d(p,q)$.
\end{definition}

\begin{theorem}[Graph criterion]\label{f:criterion}
The lattice $L$ is the full congruence lattice of a finite algebra if and only
if there is an $L$-colored complete graph $(A,d)$ satisfying all three
conditions below.
\begin{enumerate}[label=\textup{(\roman*)}]
\item\label{f:graph-triangle} If $p\ne q$, then $d(p,q)\ne0$, and
\[
 d(p,q)\le d(p,u)\vee d(u,q)\qquad(p,q,u\in A).
\]
\item\label{f:graph-reflection} Whenever $a\not\le b$ in $L$, there are vertices
$p,q$ with $d(p,q)\le a$ and $d(p,q)\not\le b$.
\item\label{f:graph-connectivity} For every nonempty edge set $S$, vertices
$p,q$ satisfying $d(p,q)\le a_S$ lie in the same connected component of
$\Gamma_S$.
\end{enumerate}
For each fixed $|A|$, these conditions involve only lattice inequalities,
finite vertex labelings, and paths.
\end{theorem}

\begin{proof}
Suppose first that $\phi:L\to\operatorname{Con}(\mathcal A)$ is an isomorphism,
where $\mathcal A$ has finite carrier $A$. Define
\[
 d(p,q)=\phi^{-1}\bigl(\operatorname{Cg}_{\mathcal A}(p,q)\bigr),
\]
where $\operatorname{Cg}_{\mathcal A}(p,q)$ is the least congruence containing
$(p,q)$. Equality is the bottom congruence, so the diagonal colors are zero
and the other colors are nonzero. The congruence
$\operatorname{Cg}_{\mathcal A}(p,u)\vee
\operatorname{Cg}_{\mathcal A}(u,q)$ contains $(p,q)$ by transitivity, proving
\ref{f:graph-triangle}. If $a\not\le b$, order reflection of $\phi$ gives a
pair in $\phi(a)\setminus\phi(b)$, and that pair proves
\ref{f:graph-reflection}.

A \emph{one-variable translation} of an operation of $\mathcal A$ is obtained
by leaving one argument free and fixing all its other arguments. Every such
translation preserves every congruence, and is therefore an admissible
self-placement. Admissible self-placements are closed under composition.
Consequently a translation sends a marked adjacency of $\Gamma_S$ to another
marked adjacency, possibly a loop. It follows, by applying the translation
to paths, that the equivalence $\rho_S$ of belonging to the same connected
component is stable under all translations. An equivalence with this
property is a congruence: change the entries of an operation's input tuple
one at a time. Nullary operations impose no further condition.

The identity placement is admissible, so $\rho_S$ contains each pair in $S$.
Its corresponding lattice element is thus at least $a_S$. If
$d(p,q)\le a_S$, then $(p,q)\in\rho_S$, which proves
\ref{f:graph-connectivity}.

Conversely, suppose $(A,d)$ satisfies the three conditions. Give $A$ one unary
operation for each admissible self-placement. There are finitely many such
maps, so this defines an allowed finite algebra $\mathcal B$. For $a\in L$ set
\[
 \theta_a=\{(p,q)\in A^2:d(p,q)\le a\}.
\]
The diagonal convention, symmetry, and \ref{f:graph-triangle} show that
$\theta_a$ is an equivalence relation. Admissibility makes it a congruence of
$\mathcal B$. The map $a\mapsto\theta_a$ is order preserving, and
\ref{f:graph-reflection} makes it order reflecting.

Let $\theta$ be any congruence of $\mathcal B$, and take for $S$ all its edges,
including all loops. This is a nonempty set, even if $A$ has one vertex.
Put $a=a_S$. By construction $\theta\subseteq\theta_a$. Each marked adjacency
of $\Gamma_S$ belongs to $\theta$, since every admissible placement is an
operation of $\mathcal B$; hence so does every connected pair. Condition
\ref{f:graph-connectivity} gives $\theta_a\subseteq\theta$. Thus
$\theta=\theta_a$, proving surjectivity and the desired isomorphism. This
argument includes the one-element lattice, represented by the one-vertex
graph and the one-element algebra.
\end{proof}

\begin{corollary}[Enumeration and bounds]\label{f:cor-witness-bound}
Representability is semidecidable. Moreover, a total algorithm deciding
representability of finite lattice order tables exists if and only if there
is a total computable function $b:\mathbb N_{>0}\to\mathbb N_{>0}$ such that
every representable lattice $L$ has a graph witness in \Cref{f:criterion} with
at most $b(|L|)$ vertices.
\end{corollary}
\begin{proof}
For each positive integer $k$, enumerate all colorings on a $k$-element vertex
set and check the three finite conditions. Running through $k=1,2,\ldots$
halts affirmatively exactly on representable lattices. A computable bound
$b$ makes this search finite and decides representability. Conversely, given
a decider, enumerate the finitely many order tables on each fixed
$n$-element set, discard tables that are not lattices, and use the decider to
retain the representable ones. Run the witness search to completion for each
retained table, and take the maximum of the resulting witness sizes and $1$.
This computes the required $b(n)$.
\end{proof}

The later undecidability argument will therefore also exclude every such
computable witness bound.
